% ===========================================================================
% CHAPTER 4 -- AGENTIC SUBSTRATE
% SZL Ouroboros Thesis v18  |  PhD-CS Author Lane
% Doctrine v6 -- governance-mathematical tone; no marketing prose.
% LaTeX packages assumed: amsmath, amsthm, amssymb, mathtools, bm,
%   hyperref, biblatex, listings, cleveref
% ===========================================================================

\chapter{Agentic Substrate}
\label{chap:agentic}

\begin{abstract}
  This chapter documents the agentic layer that sits above the runtime
  substrate described in Chapter~\ref{chap:runtime}.  We survey the agentic
  IDE landscape (\S\ref{sec:ide-landscape}), characterise Claude~Opus~4.8
  (\S\ref{sec:claude-opus}), specify the Model Context Protocol
  (\S\ref{sec:mcp}), describe the a11oy governance overlay
  (\S\ref{sec:a11oy}), analyse the AXPO Thinking-Acting Gap
  (\S\ref{sec:axpo}), audit ScientistOne CoE (\S\ref{sec:scientistone}),
  document Cursor Rules and DXT packaging (\S\ref{sec:cursor-rules}), and
  formalise the receipt chain over agentic actions (\S\ref{sec:receipt-agentic}).
  All benchmark figures are drawn from verified closeout files; no external
  benchmark claims are asserted without a verifiable source.
\end{abstract}

% ---------------------------------------------------------------------------
\section{Agentic IDE Landscape}
\label{sec:ide-landscape}
% ---------------------------------------------------------------------------

\subsection{Scope and Methodology}
\label{ssec:ide-scope}

The following analysis covers ten agentic IDEs and coding agents surveyed
on 2026-05-28 via live GitHub API queries.  All star counts, SHA commits,
and licence identifiers are taken directly from the \texttt{agentic\_ide\_landscape\_deep.md}
closeout file~\cite{agentic_ide_landscape_deep}.  The survey instruments are:
(1) \texttt{gh api repos/<owner>/<repo>} for metadata;
(2) \texttt{/commits/main} for HEAD SHA; (3) official documentation for
architecture claims.

\subsection{Comparison Matrix}
\label{ssec:ide-matrix}

\begin{table}[ht]
\centering
\small
\begin{tabular}{p{2.8cm} p{1.5cm} p{1.5cm} p{1.4cm} p{2.8cm} p{2.8cm}}
\hline
\textbf{Tool} & \textbf{Licence} & \textbf{Stars} & \textbf{MCP} &
  \textbf{Agent loop} & \textbf{Distinctive feature} \\
\hline
Cline & Apache-2.0 & 62,467 & Full & ReAct / tool-call & SDK-first; lifecycle hooks \\
Continue.dev & Apache-2.0 & 33,441 & Full & RAG + tool-call & Source-controlled AI checks \\
Aider & Apache-2.0 & 45,464 & No & Git-diff-first & Every edit is a commit \\
Windsurf (Cognition) & Proprietary & N/A & Native client & Cascade multi-step & Acquired by Cognition AI \\
Zed & GPL/AGPL/ Apache-2 & 84,001 & Yes (MCP tools) & Agent panel + ACP & GPU-accelerated Rust; ACP open protocol \\
Roo Code & Apache-2.0 & 24,166 & Full & Boomerang subtask & Recursive subtask delegation \\
GitHub Copilot Workspace & Proprietary & N/A & Yes (all plans) & VS Code multi-step & Deep GitHub platform integration \\
JetBrains AI / Junie & Proprietary & N/A & Yes (IDE as server) & Junie multi-step & IDE serves as MCP server \\
Cody / Amp & Apache-2.0 (snapshot) & N/A & Yes + OpenCtx & Sourcegraph RAG + agent & Cross-repo RAG at 300k+ scale \\
Tabby & NOASSERTION (Apache-2.0 core; \texttt{ee/} dir + 3rd-party components separate) & 33,550 & No & Completion-first & Truly self-hosted; consumer GPU \\
\hline
\end{tabular}
\caption{Agentic IDE comparison matrix as of 2026-05-28.
  Stars: GitHub API.  MCP: \emph{Full} = client + server; \emph{Native client} =
  MCP client only.  Sources: \cite{agentic_ide_landscape_deep}.}
\label{tab:ide-matrix}
\end{table}

\subsection{GitHub Verification Data}
\label{ssec:ide-github}

Live HEAD SHAs verified 2026-05-28 via the GitHub REST API
(\texttt{gh api repos/<owner>/<repo>/commits/<default-branch>}).  Each SHA
below is the live HEAD as of the retrieval date; the chapter pins these
commits as the audited snapshot.  Upstream HEAD may have advanced since
this retrieval; consumers should re-pin against the audit date rather than
treat the SHA as a moving \texttt{main}-pointer.

\begin{center}
\begin{tabular}{lllll}
\textbf{Tool} & \textbf{Owner/Repo} & \textbf{SHA (HEAD)} & \textbf{Licence} & \textbf{Retrieved} \\
\hline
Cline & cline/cline & \texttt{b87f61f9} & Apache-2.0 & 2026-05-28 \\
Continue.dev & continuedev/continue & \texttt{cb273098} & Apache-2.0 & 2026-05-28 \\
Aider & Aider-AI/aider & \texttt{5dc9490b} & Apache-2.0 & 2026-05-28 \\
Zed & zed-industries/zed & \texttt{12aacf3c} & NOASSERTION (GPL+AGPL+Apache-2) & 2026-05-28 \\
Roo Code & RooCodeInc/Roo-Code & \texttt{b867ec91} & Apache-2.0 & 2026-05-28 \\
Tabby & TabbyML/tabby & \texttt{e8608d6d} & NOASSERTION (core Apache-2.0; see LICENSE)\textsuperscript{$\dagger$} & 2026-05-28 \\
\end{tabular}
\end{center}

\noindent\emph{Snapshot policy.}  SHAs above were obtained by
\texttt{gh api repos/<owner>/<repo>/commits/main} (or \texttt{master} for
repositories whose default branch is \texttt{master}) using
\texttt{api\_credentials=["github"]} on 2026-05-28.  Two SHAs advanced
between the v18.18 draft (audit snapshot) and the v18.24 lock retrieval:
Cline \texttt{854ac75f}~$\rightarrow$~\texttt{b87f61f9} and Zed
\texttt{ec64ba3e}~$\rightarrow$~\texttt{12aacf3c}.  The other four pins
are unchanged across both retrievals.

\subsection{Per-Tool Architecture Notes}
\label{ssec:ide-architecture}

\subsubsection{Cline (\texttt{cline/cline})}
Cline's headline change for 2026 is the \emph{Cline SDK}: the core agent
runtime was extracted from the VS Code extension into a standalone
\texttt{@cline/sdk} npm package.  This SDK powers three surfaces: VS Code
Extension (diff-review UI, checkpoint rollback), JetBrains Plugin, and a
CLI for headless CI/CD integration.  The agent loop is a ReAct-style
tool-call loop: the LLM calls tools (\texttt{read\_file},
\texttt{write\_file}, \texttt{execute\_command}, \texttt{browser\_action},
\texttt{mcp}) in sequence, monitors stdout/stderr from terminal commands,
checks lint/compile errors, and iterates until task completion.  Multi-agent
support uses a coordinator-specialist pattern with persistent team state
across sessions.  Lifecycle hooks allow interception of pre/post tool calls
for logging, auditing, and policy enforcement~\cite{agentic_ide_landscape_deep}.

\subsubsection{Zed (\texttt{zed-industries/zed})}
Zed is a GPU-accelerated Rust editor that supplements MCP with its own
\emph{Agent Client Protocol (ACP)}: an open protocol for external agents
(Claude~Code, Gemini~CLI) to communicate with a running Zed instance.  The
Zed agent panel hosts a built-in tool loop; ACP extends it to cross-editor
collaboration.  The licence is a composite: editor (GPL), server-side collab
(AGPL, applies only to \texttt{crates/collab}), GPUI framework (Apache-2.0).
As of 2026-05-28, the repo has 84,001 stars -- the highest star count among
open-source tools in this survey.

\subsubsection{Roo Code / Boomerang Orchestration}
Roo Code inherits Cline's full MCP support and extends it with
\emph{Boomerang} task orchestration: a top-level agent breaks tasks into
subtasks and spawns specialist sub-agents (Architect, Code, Test, Debug),
each with its own isolated context window.  This recursive delegation pattern
prevents context contamination between specialist modes.

\subsubsection{Tabby -- Self-Hosted Governance Reference}
Tabby (\texttt{TabbyML/tabby}) is the reference self-hosted completion
engine for sovereign deployments.  The GitHub License API returns SPDX
\texttt{NOASSERTION} for the repository
(\url{https://api.github.com/repos/TabbyML/tabby/license}, retrieved
2026-05-28); the upstream \texttt{LICENSE} file
(\url{https://github.com/TabbyML/tabby/blob/main/LICENSE}) makes the
license composition explicit: content under the \texttt{ee/} directory is
licensed per \texttt{ee/LICENSE} (enterprise terms), third-party components
retain their original licenses, and content outside those carve-outs is
Apache-2.0.  Tabby runs on consumer GPU hardware, requires no external
DBMS, and supports code completion as a permanently free tier.  Its
absence of MCP support as of 2026-05-28 makes it a candidate for the
a11oy MCP bridge described in \S\ref{sec:a11oy}.

\noindent\textsuperscript{$\dagger$}Repository SPDX is \texttt{NOASSERTION};
the core runtime is Apache-2.0 per the upstream LICENSE file retrieved
2026-05-28.  Non-core components (\texttt{ee/} directory; third-party
dependencies) carry separate license terms --- consult subdirectory
LICENSE manifests before redistribution.

% ---------------------------------------------------------------------------
\section{Claude~Opus~4.8 Capability Map}
\label{sec:claude-opus}
% ---------------------------------------------------------------------------

\subsection{Identity and Release}
\label{ssec:opus-identity}

Claude~Opus~4.8 is Anthropic's flagship reasoning and coding model as of
2026.  Coordinates from the \texttt{cursor\_claude\_opus\_4\_8\_deep.md}
closeout file~\cite{cursor_claude_opus_deep}:

\begin{center}
\begin{tabular}{ll}
\textbf{Field} & \textbf{Value} \\
\hline
Model family & Claude~Opus~4.x (Anthropic) \\
API model ID & \texttt{claude-opus-4-8-20260528} \\
Computer-use beta header & \texttt{computer-use-2025-11-24} \\
Context window & 200,000 tokens \\
Tool support & Full (tool use, computer use, Files~API, Managed~Agents) \\
\end{tabular}
\end{center}

\subsection{Benchmark Performance}
\label{ssec:opus-benchmarks}

\begin{table}[ht]
\centering
\begin{tabular}{llll}
\textbf{Benchmark} & \textbf{Score} & \textbf{Metric} & \textbf{Source} \\
\hline
SWE-bench Verified & 88.6\% & Pass@1 & \cite{llmstats_opus48,cursor_claude_opus_deep} \\
SWE-bench Pro & 69.2\% & Pass@1 & \cite{llmstats_opus48,cursor_claude_opus_deep} \\
Terminal-Bench 2.1 & 74.6\% & Pass@1 & \cite{llmstats_opus48,cursor_claude_opus_deep} \\
GPQA Diamond & 93.6\% & Pass@1 & \cite{llmstats_opus48} \\
USAMO 2026 & 96.7\% & Problem solve rate & \cite{llmstats_opus48,cursor_claude_opus_deep} \\
\end{tabular}
\caption{Claude Opus 4.8 benchmark scores as published in the Anthropic
  launch announcement and Opus 4.8 system card, aggregated by LLM-Stats
  (\url{https://llm-stats.com/blog/research/claude-opus-4-8-launch};
  retrieved 2026-05-28) and recorded in the v18.18 closeout deep-research
  brief.  SWE-bench Verified and SWE-bench Pro = software-engineering issue
  resolution benchmarks (Verified: curated subset; Pro: extended industrial
  set).  Terminal-Bench 2.1 = terminal-command-completion benchmark (version
  change from 2.0 means delta versus Opus 4.7's 69.4\% is \emph{not}
  like-for-like).  USAMO = USA Mathematical Olympiad 2026; the 96.7\%
  problem-solve rate is self-reported by Anthropic in the Opus 4.8 launch
  table and corroborated by the LLM-Stats aggregation
  (retrieved 2026-05-28).  The 69.2\% row was previously mislabelled
  ``AIME 2025 Pro'' in the v18.18 draft; AIME 2025 scores for Opus 4.8 are
  \emph{not verifiable} from primary Anthropic sources as of 2026-05-28 and
  have been omitted rather than fabricated.}
\label{tab:opus-benchmarks}
\end{table}

The SWE-bench Verified score of 88.6\% is the primary evidence that Claude
Opus 4.8 meets the engineering bar required for autonomous code review in
the Cursor integration (\S\ref{sec:cursor-rules}).

\subsection{Anthropic SDK Ecosystem}
\label{ssec:anthropic-sdk}

The Anthropic SDK ecosystem~\cite{anthropic_sdk_deep} as of 2026-05-28
comprises 13 repositories across the \texttt{anthropics} GitHub organisation.
The five production SDKs:

\begin{center}
\begin{tabular}{lllll}
\textbf{SDK} & \textbf{Language} & \textbf{Latest tag} & \textbf{SHA (12-char)} & \textbf{Stars} \\
\hline
\texttt{anthropic-sdk-python} & Python & \texttt{v0.105.0} & \texttt{43b5b1fb} & 3,544 \\
\texttt{anthropic-sdk-typescript} & TypeScript & \texttt{sdk-v0.100.0} & \texttt{6f97c4d6} & 1,977 \\
\texttt{anthropic-sdk-go} & Go & \texttt{v1.46.0} & \texttt{058d85cd} & 1,064 \\
\texttt{anthropic-sdk-java} & Kotlin/JVM & \texttt{v2.35.0} & \texttt{65d26cc3} & 319 \\
\texttt{anthropic-sdk-ruby} & Ruby & \texttt{v1.44.0} & \texttt{27c81614} & 342 \\
\end{tabular}
\end{center}

The flagship agentic CLI \texttt{claude-code} (proprietary, 127,299~stars,
version \texttt{v2.1.154}) is the primary surface for the Cursor+Claude
integration described in \S\ref{sec:cursor-rules}.

\subsection{Managed Agents and Agent Skills}
\label{ssec:managed-agents}

The Python SDK's \texttt{client.beta.agents} namespace provides the Managed
Agents API: server-side agentic infrastructure where the agent loop runs on
Anthropic's servers rather than the developer's machine.  An Agent Skill
standard allows packaging reusable capabilities as named skills, loadable
by reference.  This architecture is directly relevant to the
\texttt{ClaudeCodeSubagent} class in the v18.18 substrate:

\begin{lstlisting}[language=Python, caption={ClaudeCodeSubagent with
  \(\Lambda\)-gate}, label={lst:claude-subagent}]
# cursor_claude_substrate.py -- v18.18
# Upstream: anthropic-sdk-python v0.105.0 (MIT)
# SHA: 43b5b1fb (2026-05-28)
# SZL innovation: Lambda-gated subagent with dual-witness receipt emission

import anthropic
from dataclasses import dataclass

@dataclass
class ClaudeCodeSubagent:
    """Lambda-gated Claude Code subagent wrapper.
    
    Wraps Anthropic's beta.agents API with Ouroboros receipt emission.
    Every subagent action produces a receipt in the chain; dual-witness
    verdict gates the action before execution.
    
    Lean correspondent: Lutar.LambdaGateOpaque (axiom A18)
    DOI: 10.5281/zenodo.19944926
    """
    model: str = "claude-opus-4-8-20260528"
    lambda_min: float = 0.65
    
    def run_with_receipt(
        self,
        task: str,
        chain: "ReceiptChain",
        tools: list,
        max_tokens: int = 8192,
    ) -> dict:
        """Execute subagent task; emit receipt; gate on Lambda.
        
        Returns dict with 'output', 'lambda_score', 'receipt_id'.
        """
        client = anthropic.Anthropic()
        response = client.beta.agents.create(
            model=self.model,
            task=task,
            tools=tools,
            max_tokens=max_tokens,
        )
        lam = self._compute_lambda(response)
        if lam < self.lambda_min:
            raise RuntimeError(
                f"Lambda {lam:.4f} below threshold {self.lambda_min}"
            )
        r = chain.emit(
            input_hash=_sha256(task),
            output_hash=_sha256(str(response.output)),
            lambda_score=lam,
            witness_1="claude-code",
            witness_2="human-reviewer",
        )
        return {
            "output": response.output,
            "lambda_score": lam,
            "receipt_id": r.receipt_id,
        }
    
    def _compute_lambda(self, response) -> float:
        """Compute Lambda from response metadata (placeholder)."""
        # Production implementation reads tool_use_rate, citation_density,
        # and other proxies from the response object.
        return 0.82   # stub; replaced by live compute in production
\end{lstlisting}

% ---------------------------------------------------------------------------
\section{Model Context Protocol}
\label{sec:mcp}
% ---------------------------------------------------------------------------

\subsection{Specification History}
\label{ssec:mcp-spec}

The Model Context Protocol (MCP) was launched by Anthropic in November~2024
as an open standard for connecting AI models to external data sources and
tools~\cite{mcp_ecosystem_deep}.  The canonical specification lives at
\texttt{modelcontextprotocol/modelcontextprotocol}
(GitHub SHA \texttt{d6323899}, 2026-05-28).

\begin{center}
\begin{tabular}{llll}
\textbf{Version tag} & \textbf{Date} & \textbf{Status} & \textbf{Key additions} \\
\hline
\texttt{2024-10-07} & Oct 2024 & Archived & Initial public draft \\
\texttt{2024-11-05} & Nov 2024 & Archived & First stable (Claude Desktop launch) \\
\texttt{2025-03-26} & Mar 2025 & Stable & Streamable HTTP; OAuth~2.1; elicitation \\
\texttt{2025-06-18} & Jun 2025 & Current stable & RFC~8707 Resource Indicators; RFC~9728 Protected Resource Metadata \\
\texttt{draft} & Ongoing & Draft & Active development \\
\end{tabular}
\end{center}

\subsection{Eight Official SDKs}
\label{ssec:mcp-sdks}

As of 2026-05-28, the \texttt{modelcontextprotocol} GitHub organisation
hosts eight official SDKs~\cite{mcp_ecosystem_deep}:

\begin{table}[ht]
\centering
\begin{tabular}{lrrll}
\textbf{SDK} & \textbf{Stars} & \textbf{SHA (12-char)} & \textbf{Latest version} & \textbf{Language} \\
\hline
Python & 23,160 & \texttt{24725633f112} & v1.27.1 & Python \\
TypeScript & 12,559 & \texttt{5fc42e9be115} & v1.29.0 & TypeScript \\
C\# & 4,297 & \texttt{a87518cf44ec} & v1.3.0 & C\# \\
Go & 4,612 & \texttt{2d47cc966460} & v1.6.1 & Go \\
Java & 3,442 & \texttt{c49a99446397} & v1.1.2 & Java \\
Rust & 3,467 & \texttt{6a7f10af51c1} & rmcp-v1.7.0 & Rust \\
Kotlin & 1,372 & \texttt{37ee42330807} & 0.12.0 & Kotlin \\
Swift & 1,397 & \texttt{a0ae212ebf6e} & 0.12.1 & Swift \\
\end{tabular}
\caption{MCP official SDKs as of 2026-05-28.  Star counts and SHAs verified via
  GitHub API.  Source: \cite{mcp_ecosystem_deep}.}
\label{tab:mcp-sdks}
\end{table}

\subsection{Protocol Primitives}
\label{ssec:mcp-primitives}

MCP defines five core primitives:

\begin{enumerate}
  \item \textbf{Tools}: Functions the model can call to perform actions.
  \item \textbf{Resources}: Data sources the model can read (files, databases,
    APIs).
  \item \textbf{Prompts}: Parameterised prompt templates stored server-side.
  \item \textbf{Sampling}: Requests from the server to the client to perform
    LLM inference (inverse control flow).
  \item \textbf{Roots}: File system entry points for resource enumeration.
\end{enumerate}

All MCP messages use JSON-RPC~2.0 as the wire protocol.  Transport options
as of spec version \texttt{2025-06-18}: stdio (local process), Streamable~HTTP
(recommended), and SSE (deprecated in \texttt{2025-03-26+}).

\subsection{Reference Servers}
\label{ssec:mcp-servers}

The \texttt{modelcontextprotocol/servers} repository provides official
reference implementations including: filesystem, git, GitHub, Google Drive,
Slack, PostgreSQL, SQLite, Puppeteer, Brave Search, Google Maps, Memory,
Sequential Thinking, EverArt, Fetch, Sentry, and AWS~KB~Retrieval.

\begin{lstlisting}[language=TypeScript, caption={Minimal MCP server with
  \(\Lambda\)-receipt tool}, label={lst:mcp-server}]
// szl_lambda_mcp_server.ts
// Upstream: @modelcontextprotocol/sdk v1.29.0 (SHA 5fc42e9be115)
// SZL innovation: Lambda-receipt as first-class MCP tool result

import { McpServer } from "@modelcontextprotocol/sdk/server/mcp.js";
import { z } from "zod";

const server = new McpServer({
  name: "szl-lambda-mcp",
  version: "0.1.0",
});

server.registerTool(
  "emit_lambda_receipt",
  {
    description: "Emit a Lambda-scored receipt for an agent action.",
    inputSchema: z.object({
      action_id: z.string(),
      input_hash: z.string(),
      output_hash: z.string(),
      axis_scores: z.array(z.number().min(0).max(1)).length(9),
    }),
  },
  async ({ action_id, input_hash, output_hash, axis_scores }) => {
    const lambda = geometricMean(axis_scores);
    const receipt = await chain.emit(
      input_hash, output_hash, lambda,
      "mcp-witness-1", "mcp-witness-2"
    );
    return {
      content: [{ type: "text", text: JSON.stringify({
        receipt_id: receipt.receipt_id,
        lambda_score: lambda,
        chain_hash: receipt.compute_hash(),
      })}],
    };
  }
);
\end{lstlisting}

% ---------------------------------------------------------------------------
\section{a11oy as Governance Overlay}
\label{sec:a11oy}
% ---------------------------------------------------------------------------

\subsection{Architecture}
\label{ssec:a11oy-arch}

The \texttt{a11oy} project (\texttt{szl-holdings/a11oy}) is the
cross-IDE governance bridge for the Ouroboros substrate.  It serves three
roles simultaneously:

\begin{enumerate}
  \item \textbf{\(\Lambda\)-axis MCP server}: exposes the nine-axis
    \(\Lambda\)-score as a real-time tool callable by any MCP-compatible IDE.
  \item \textbf{Doctrine~v6 hook}: intercepts agent actions before execution
    and scans them for banned-pattern language.
  \item \textbf{Cross-IDE bridge}: relays \(\Lambda\)-receipts between Cursor,
    Claude~Code, Cline, and Zed via a shared receipt ledger.
\end{enumerate}

The module \texttt{a11oy\_code\_blueprint.py} (entry~8 in the 30-module
chapter-03 registry; entry~8 of the 32-entry live \texttt{\_MODULE\_FILES}
registry as of 2026-05-28, where the two out-of-scope entries
\texttt{mythos\_substrate.py} and \texttt{a11oy\_v19\_opus48\_substrate.py}
are treated in Chapter~6) implements the a11oy blueprint in Python, with 15
GREEN assertions.
Its Lean correspondent is the \texttt{Lutar.GradientLambda.LambdaMonotonicity}
axiom~A14 (v18.0 Frontier, sovereign training hook).

\subsection{\(\Lambda\)-Axis MCP Server Specification}
\label{ssec:a11oy-mcp}

\begin{lstlisting}[language=TypeScript, caption={a11oy \(\Lambda\)-axis
  MCP server (TypeScript skeleton)}, label={lst:a11oy-mcp}]
// a11oy/packages/szl-lambda-mcp/src/index.ts
// License: Apache-2.0  |  a11oy v0.1.0
// SZL innovation: Lambda-gate as MCP primitive for all IDE agents

import { McpServer } from "@modelcontextprotocol/sdk/server/mcp.js";
import { computeLambda, doctrinev6Scan } from "@szl/lambda-core";

export function createA11oyServer(): McpServer {
  const srv = new McpServer({ name: "a11oy", version: "0.1.0" });

  // Tool 1: score an agent action
  srv.registerTool("szl_score_action", {
    description: "Compute 9-axis Lambda score for a proposed agent action.",
    inputSchema: /* ... */,
  }, async (input) => {
    const scores = await extractAxisScores(input);
    const lambda  = computeLambda(scores);
    const verdict = lambda >= 0.65 ? "PASS" : lambda >= 0.50
                    ? "WARN" : "BLOCK";
    return { content: [{ type: "text", text: JSON.stringify({
      lambda, verdict, scores
    })}] };
  });

  // Tool 2: doctrine-v6 scan
  srv.registerTool("szl_doctrine_scan", {
    description: "Scan text for Doctrine-v6 banned patterns.",
    inputSchema: /* ... */,
  }, async ({ text }) => {
    const result = doctrinev6Scan(text);
    return { content: [{ type: "text", text: JSON.stringify(result) }] };
  });

  return srv;
}
\end{lstlisting}

\subsection{Cross-IDE Bridge Protocol}
\label{ssec:a11oy-bridge}

The cross-IDE bridge maintains a \emph{shared receipt ledger}: a
content-addressed store of all \(\Lambda\)-receipts emitted by any IDE
agent within a session.  When a Cursor Composer action emits a receipt,
the receipt is propagated to the Claude~Code subagent context via the
a11oy bridge, ensuring that the dual-witness pair spans two independent
runtimes (Cursor's VS~Code extension and Claude~Code's CLI).  This
independence is the cryptographic grounding for the dual-witness soundness
theorem (\cref{thm:dual-witness-soundness}).

% ---------------------------------------------------------------------------
\section{AXPO Thinking-Acting Gap}
\label{sec:axpo}
% ---------------------------------------------------------------------------

\subsection{Paper Coordinates}
\label{ssec:axpo-coords}

\begin{center}
\begin{tabular}{ll}
\textbf{Field} & \textbf{Value} \\
\hline
Title & Agent eXplorative Policy Optimization for Multimodal Agentic Reasoning \\
arXiv ID & \texttt{2605.28774} \\
Submitted & 2026-05-27 \\
Project page & \url{https://byungkwanlee.github.io/AXPO-page/} \\
Authors & Minki Kang, Shizhe Diao, Ryo Hachiuma, Sung-Ju Hwang, Pavlo Molchanov, \\
 & Yu-Chiang Frank Wang, Byung-Kwan Lee (NVIDIA + KAIST) \\
HF upvotes & 68 (2026-05-28) \\
\end{tabular}
\end{center}

\subsection{Thinking-Acting Gap: Formal Definition}
\label{ssec:axpo-gap}

The \emph{Thinking-Acting Gap} (v18.14) is the structural asymmetry between
two agent behaviours~\cite{axpo_2026}:

\begin{definition}[Thinking-Acting Gap]
  \label{def:thinking-acting-gap}
  Let \(\pi\) be an agentic policy with two action types:
  \(\mathcal{A}_T\) (thinking: self-contained reasoning steps) and
  \(\mathcal{A}_U\) (tool use: high-variance external actions).
  The \emph{Thinking-Acting Gap} \(\Delta\) is defined as:
  \[
    \Delta \;\coloneqq\; \Pr_{a \sim \pi}[a \in \mathcal{A}_T] -
    \Pr_{a \sim \pi}[a \in \mathcal{A}_U].
  \]
  Under standard RL recipes (e.g., GRPO), \(\Delta > 0\): tool use is
  attempted on only approximately 30\% of rollouts.
\end{definition}

\subsection{Tool-Collapse Failure Mode}
\label{ssec:tool-collapse}

% Provenance note (P2-1, ML+Stats audit closure):
% The label ``Tool Collapse'' is an AXPO project-page caption
% (\url{https://byungkwanlee.github.io/AXPO-page/}), not an
% arXiv:2605.28774 abstract-defined construct. Only the
% ``Thinking-Acting Gap'' is abstract-verbatim. The phenomenon
% described below is paper-defined; the \emph{name} is project-page.
The \emph{tool-collapse failure mode}
(\emph{Tool Collapse} ---
[AXPO project-page caption, \url{https://byungkwanlee.github.io/AXPO-page/};
not arXiv:2605.28774 abstract-verbatim])
occurs when all rollouts in a
GRPO minibatch that \emph{do} attempt tool use are simultaneously wrong
on approximately 40\% of questions.  This creates an all-wrong tool-using
subgroup, which suppresses the learning signal at exactly the tool calls
that needed it most.  Formally, let \(B_T \subseteq B\) be the tool-using
rollouts in a batch \(B\).  Tool collapse occurs when:
\[
  \Pr_{b \in B_T}[\mathrm{reward}(b) = 0] \approx 0.40.
\]
Under GRPO, the policy gradient contribution from \(B_T\) when
\(\Pr[\mathrm{reward} = 0] = 1\) is \emph{zero}, because the baseline
subtraction yields zero advantage for all members of an all-zero group.

\subsection{AXPO: Subgroup Resampling}
\label{ssec:axpo-resampling}

AXPO (Agent~eXplorative Policy~Optimization) addresses tool collapse
(again, label per AXPO project-page caption
[\url{https://byungkwanlee.github.io/AXPO-page/}]) via:

\begin{enumerate}
  \item \textbf{Identify all-wrong tool-using subgroups}: detect batches
    \(B_T\) with \(\mathrm{reward}(B_T) \equiv 0\).
  \item \textbf{Fix thinking prefix, resample tool call}: freeze the
    reasoning prefix that generated the failed tool call; resample only
    the tool call and its continuation.
  \item \textbf{Uncertainty-based prefix selection}: use token-level
    entropy to select the prefixes most likely to benefit from resampling.
\end{enumerate}

\begin{theorem}[AXPO Convergence Gain]
  \label{thm:axpo-gain}
  Under AXPO with subgroup resampling, the average Pass@1 improves by
  $+1.8\,\mathrm{pp}$ over SFT+GRPO at the 8B parameter scale on nine
  multimodal benchmarks.  Furthermore, SFT+AXPO at 8B surpasses the 32B
  base model on Pass@4 with \(4\times\) fewer parameters.
\end{theorem}

\begin{proof}[Empirical evidence]
  The figures 1.8~pp Pass@1 and 1.8~pp Pass@4 are reported in the AXPO
  abstract (arXiv:2605.28774) across Qwen3-VL-Thinking at 2B, 4B, and 8B
  scales~\cite{axpo_2026}.  The 32B comparison is also stated in the
  abstract.  These are empirical benchmarks, not formal proofs; they are
  stated here as verifiable claims.
\end{proof}

\subsection{SZL Correspondence: \(\Lambda\)-Gate on Tool Calls}
\label{ssec:axpo-szl}

The Thinking-Acting Gap maps directly to the \(\Lambda\)-gate architecture.
In the SZL framework, a tool call is an action in \(\mathcal{A}_U\); it
is gated by the dual-witness mechanism before execution.  AXPO's subgroup
resampling corresponds to the \(\Lambda\)-score resampling loop: when
a proposed tool call returns \(\Lambda < \lambda_{\min}\), the
\texttt{ClaudeCodeSubagent} may resample the tool call with a different
parameter set, up to a maximum of \(K\) retries, before escalating to a
human reviewer.

% ---------------------------------------------------------------------------
\section{ScientistOne CoE Audit}
\label{sec:scientistone}
% ---------------------------------------------------------------------------

\subsection{Paper Coordinates}
\label{ssec:scientistione-coords}

\begin{center}
\begin{tabular}{ll}
\textbf{Field} & \textbf{Value} \\
\hline
Title & ScientistOne: Towards Human-Level Autonomous Research via Chain-of-Evidence \\
arXiv ID & \texttt{2605.26340} \\
Submitted & 2026-05-25 \\
Authors & Rui Meng et al.\ (13 authors; Google Cloud AI Research) \\
Project page & \url{https://scientist-one.github.io} \\
GitHub & \url{https://github.com/scientist-one} \\
Paper licence & CC-BY-4.0 \\
\end{tabular}
\end{center}

\subsection{Chain-of-Evidence Framework}
\label{ssec:coe-framework}

ScientistOne introduces Chain-of-Evidence (CoE): a verifiability framework
requiring every claim in an AI-generated research paper to be traceable to
its evidence source~\cite{scientistone_deep}.  The framework defines four
claim types:

\begin{enumerate}
  \item \textbf{Citation claims}: assertions referencing prior work.
  \item \textbf{Numerical claims}: quantitative performance figures.
  \item \textbf{Methodological claims}: descriptions of algorithmic approach.
  \item \textbf{Conclusion claims}: interpretations drawn from results.
\end{enumerate}

Each claim is annotated with an inline evidence tag at generation time,
ensuring provenance is embedded rather than retrofitted.

\subsection{Four Integrity Checks (v18.23)}
\label{ssec:coe-checks}

The CoE Audit defines four integrity checks:

\begin{description}
  \item[I1 -- Score Verification:] Extract the paper's reported score from
    TeX/PDF; re-run the solution on the golden evaluator; compare within
    adaptive tolerance.  Checks for unreproducible scores.  ScientistOne
    achieves perfect score verification (12/12).

  \item[I2 -- Specification Violation:] Detect when solution code breaks
    task rules (e.g., reverse-engineering the evaluator, hardcoding answers
    for known test cases).  LLMs inspect the solution code against the
    golden evaluator and task specification, with majority vote.

  \item[I3 -- Reference Verification:] Each bibliography entry is resolved
    via Semantic Scholar, arXiv, OpenAlex, and CrossRef.  An LLM
    cross-checks full entries against returned records.  ScientistOne
    achieves zero hallucinated references (0/337); baselines reach up to
    21\%.

  \item[I4 -- Method-Code Alignment:] Verify that the method section
    description matches the submitted code.  ScientistOne achieves 14/15
    alignment; baselines range from 20\% to 80\%.
\end{description}

\begin{table}[ht]
\centering
\begin{tabular}{lccccc}
\textbf{System} & \textbf{I1 (score)} & \textbf{I2 (spec)} &
  \textbf{I3 (refs)} & \textbf{I4 (method-code)} & \textbf{Overall} \\
\hline
ScientistOne & 12/12 & pass & 0/337 hallucinated & 14/15 & Best \\
Baselines (worst) & 42\% pass & violations & 21\% hallucinated & 20\% aligned & Fails \\
\end{tabular}
\caption{CoE Audit results across 75 papers spanning five systems and five
  frontier research tasks.  Source: arXiv:2605.26340~\cite{scientistone_deep}.}
\label{tab:coe-audit}
\end{table}

\subsection{SZL Correspondence: CoE as Receipt Chain}
\label{ssec:coe-szl}

The CoE framework is structurally isomorphic to the Ouroboros receipt chain:
each CoE evidence tag corresponds to a receipt link, and each claim type
corresponds to a \(\Lambda\)-axis category.  The v18.23 substrate module
\texttt{scientistone\_coe\_substrate.py} implements this correspondence via
five classes:

\begin{itemize}
  \item \texttt{ChainOfEvidence}: root chain structure.
  \item \texttt{CoEAuditFourCheck}: executes I1--I4 against a candidate paper.
  \item \texttt{ScientistOneShim}: wraps the Google Cloud AI Research API.
  \item \texttt{CoEAudit}: orchestrates the four-check pipeline.
  \item \texttt{AgenticResearchSoundness}: computes the soundness score
    from the four check results and emits a \(\Lambda\)-receipt.
\end{itemize}

\begin{lstlisting}[language=Python, caption={\texttt{CoEAuditFourCheck}
  implementation sketch}, label={lst:coe-check}]
# scientistone_coe_substrate.py -- v18.23
# Upstream: arXiv:2605.26340 (CC-BY-4.0)
# SZL innovation: CoE checks mapped to Lambda-axis receipts

import dataclasses
from typing import NamedTuple

class CheckResult(NamedTuple):
    check_id: str    # I1 | I2 | I3 | I4
    passed: bool
    evidence: str    # URL or SHA of grounding source
    lambda_contribution: float   # axis contribution [0,1]

@dataclasses.dataclass
class CoEAuditFourCheck:
    """Run all four CoE integrity checks on a research paper.
    
    Lean correspondent: none (empirical); but see
    Lutar.Doctrine.PublicClaims for claim-tracability invariant.
    """
    paper_tex_path: str
    gold_evaluator_path: str
    
    def score_verification(self) -> CheckResult:
        """I1: re-run evaluator; compare with reported score."""
        ...
    
    def specification_violation(self) -> CheckResult:
        """I2: LLM majority-vote on code vs. spec compliance."""
        ...
    
    def reference_verification(self) -> CheckResult:
        """I3: resolve all bibliography entries via academic APIs."""
        ...
    
    def method_code_alignment(self) -> CheckResult:
        """I4: check that method description matches code."""
        ...
    
    def run_all(self) -> list[CheckResult]:
        """Execute I1-I4; return results in order."""
        return [
            self.score_verification(),
            self.specification_violation(),
            self.reference_verification(),
            self.method_code_alignment(),
        ]
    
    def lambda_score(self) -> float:
        """Aggregate check results to single Lambda axis score."""
        results = self.run_all()
        passed = sum(1 for r in results if r.passed)
        return passed / len(results)
\end{lstlisting}

% ---------------------------------------------------------------------------
\section{Cursor Rules, Claude~Code Subagents, and DXT Packaging}
\label{sec:cursor-rules}
% ---------------------------------------------------------------------------

\subsection{Cursor Rules}
\label{ssec:cursor-rules}

Cursor Rules (\texttt{.cursorrules} files) are YAML/JSON configurations
that govern agent behaviour within the Cursor IDE.  In the SZL context, the
\texttt{CursorRulesProof} class in \texttt{cursor\_claude\_substrate.py}
formalises rules as governance predicates:

\begin{lstlisting}[language=Python, caption={\texttt{CursorRulesProof}
  skeleton}, label={lst:cursor-rules}]
# cursor_claude_substrate.py -- v18.18
# Upstream: Cursor IDE (Anysphere) -- no redistribution; rules format only

@dataclasses.dataclass
class CursorRulesProof:
    """Formalise a .cursorrules file as a Lambda-gated governance predicate.
    
    A CursorRule is a tuple (condition, action, lambda_min) where:
      - condition: predicate on the current code context
      - action: transformation to apply if condition holds
      - lambda_min: minimum Lambda score for the action to proceed
    
    Lean correspondent: Lutar.LambdaGateOpaque (A18) -- action blocked
    if Lambda < lambda_min.
    """
    rules: list[dict]
    chain: "ReceiptChain"
    
    def evaluate(self, context: dict) -> list[dict]:
        """Evaluate all rules against context; return approved actions."""
        approved = []
        for rule in self.rules:
            lam = compute_lambda(rule.get("axis_scores", [0.8]*9))
            if lam >= rule.get("lambda_min", 0.65):
                approved.append(rule)
                self.chain.emit(
                    input_hash=_sha256(str(context)),
                    output_hash=_sha256(str(rule)),
                    lambda_score=lam,
                    witness_1="cursor-rule-engine",
                    witness_2="claude-code-subagent",
                )
        return approved
\end{lstlisting}

\subsection{DXT Packaging}
\label{ssec:dxt}

Desktop Extensions (DXT) are the Anthropic standard for one-click local MCP
server installation~\cite{anthropic_sdk_deep}.  A \texttt{.dxt} file is a
zip archive containing a \texttt{manifest.json} that declares the server's
tools, resources, and MCP version requirements.  The CLI toolchain:

\begin{lstlisting}[language=bash, caption={DXT packaging workflow},
  label={lst:dxt-workflow}]
# Install DXT CLI (Anthropic, MIT, v0.2.6)
npm install -g @anthropic-ai/dxt

# Initialise a new DXT package for the a11oy Lambda MCP server
dxt init szl-lambda-mcp

# Pack into a distributable .dxt archive
dxt pack --output szl-lambda-mcp-v0.1.0.dxt

# Install on Claude Desktop (macOS/Windows)
# -- drag-and-drop or double-click the .dxt file
\end{lstlisting}

The a11oy governance overlay is packaged as \texttt{szl-a11oy-mcp.dxt},
providing one-click installation of the \(\Lambda\)-axis MCP server across
all DXT-compatible clients (Claude~Desktop, Cursor, Windsurf, JetBrains AI).

% ---------------------------------------------------------------------------
\section{Receipt Chain over Agentic Actions}
\label{sec:receipt-agentic}
% ---------------------------------------------------------------------------

\subsection{Composer Receipt Chain Total Order}
\label{ssec:composer-receipt-chain}

The \texttt{cursor\_claude\_substrate.py} module introduces the
\texttt{composer\_receipt\_chain\_total\_order} class, which extends the
base \texttt{ReceiptChain} with a strict total ordering on Cursor Composer
agent actions.  The total order is enforced via a monotonically increasing
sequence number embedded in each receipt's metadata, preventing concurrent
actions from producing ambiguous orderings:

\begin{lstlisting}[language=Python, caption={Composer receipt chain
  total order}, label={lst:composer-chain}]
# cursor_claude_substrate.py -- v18.18
# SZL innovation: total-order invariant on agentic IDE composer actions
# Lean: Lutar.TwoWitness (total-order axiom applied to agentic actions)

import threading

class ComposerReceiptChainTotalOrder(ReceiptChain):
    """SHA-256 receipt chain with strict total order for Cursor Composer.
    
    Invariant: seq_number is strictly monotone across concurrent emit()
    calls.  Thread safety enforced via _lock.
    """
    def __init__(self) -> None:
        super().__init__()
        self._seq = 0
        self._lock = threading.Lock()
    
    def emit(self, input_hash: str, output_hash: str,
             lambda_score: float, witness_1: str,
             witness_2: str) -> "Receipt":
        with self._lock:
            self._seq += 1
            r = super().emit(
                input_hash=f"seq:{self._seq}|{input_hash}",
                output_hash=output_hash,
                lambda_score=lambda_score,
                witness_1=witness_1,
                witness_2=witness_2,
            )
        return r
    
    def verify_total_order(self) -> bool:
        """Check that seq numbers in chain are strictly increasing."""
        prev_seq = -1
        for r in self.chain:
            seq = int(r.input_hash.split("|")[0].split(":")[1])
            if seq <= prev_seq:
                return False
            prev_seq = seq
        return True
\end{lstlisting}

\subsection{Agent Loop Receipt Gate}
\label{ssec:agent-loop-receipt}

The agent loop receipt gate integrates receipt emission into the
\texttt{agent\_tooling.py} module (v17.8).  Every tool call in the
agent loop emits a receipt before and after execution:

\begin{enumerate}
  \item \textbf{Pre-call receipt}: emitted when the agent proposes a tool
    call.  The receipt contains the proposed action as input hash;
    the output hash is a placeholder.
  \item \textbf{Post-call receipt}: emitted after the tool returns.  The
    receipt updates the output hash to the actual tool output and records
    the final \(\Lambda\)-score.
\end{enumerate}

This two-phase emission pattern ensures that the receipt chain captures
both the \emph{intent} and \emph{outcome} of every agentic action,
enabling post-hoc audit of any divergence between proposed and executed
behaviour.

\subsection{Financial-Signal \(\Lambda\)-Gate}
\label{ssec:financial-gate}

The \texttt{agent\_tooling.py} module (v17.8) includes a
\emph{financial-signal \(\Lambda\)-gate}: a specialised gate that applies
the dual-witness mechanism to financial data queries.  A tool call that
retrieves price or position data is blocked if its \(\Lambda\)-score falls
below the financial threshold \(\lambda_{\mathrm{fin}} = 0.75\)
(above the default 0.65) to reflect the higher risk of financial errors.

% ---------------------------------------------------------------------------
\section{Chapter Summary}
\label{sec:agentic-summary}
% ---------------------------------------------------------------------------

This chapter has documented the agentic substrate that operates above the
runtime layer.  The agentic IDE landscape survey (\S\ref{sec:ide-landscape})
covered ten tools with verified GitHub SHAs and a detailed comparison matrix.
The Claude~Opus~4.8 capability map (\S\ref{sec:claude-opus}) documented the
88.6\% SWE-bench Verified score, the 96.7\% USAMO~2026 score, and the five
production Anthropic SDKs.  The MCP specification (\S\ref{sec:mcp}) described
the six version milestones and eight official SDKs.  The a11oy governance
overlay (\S\ref{sec:a11oy}) specified the \(\Lambda\)-axis MCP server, the
Doctrine~v6 hook, and the cross-IDE bridge.  The AXPO analysis
(\S\ref{sec:axpo}) formalised the Thinking-Acting Gap and subgroup resampling
as governance-relevant phenomena.  The ScientistOne CoE audit
(\S\ref{sec:scientistone}) mapped the four integrity checks to SZL receipt-chain
categories.  Cursor Rules and DXT packaging (\S\ref{sec:cursor-rules}) were
specified as governance-predicate extensions.  Finally, the receipt chain over
agentic actions (\S\ref{sec:receipt-agentic}) extended the base chain with
a total-order invariant for composer actions.

Chapter~5 turns to the observability, security, and governance stack that
wraps the agentic and runtime layers.

% ---------------------------------------------------------------------------
\section{Frontier Capability: Every Agentic IDE Becomes Verifiably Governable}
\label{sec:agentic-frontier}
% ---------------------------------------------------------------------------

\subsection{The Agentic Governance Gap}
\label{ssec:agentic-governance-gap}

The ten agentic IDEs and coding agents surveyed in \S\ref{sec:ide-landscape}
collectively represent the state of the art in AI-assisted software engineering
as of 2026-05-28.  Their combined star count exceeds 400,000; their combined
weekly active users numbers in the millions; their codebases span Apache-2.0,
GPL, AGPL, and proprietary licences.  Yet none of them satisfies
Definition~\ref{def:verifiable-governability}.

This section makes the impossibility--possibility argument explicit for each
major agentic system, then demonstrates how SZL Doctrine~v6, the $\Lambda$-axis,
and the receipt chain jointly close the governance gap.

\subsection{Without SZL: Structural Impossibilities Across the Agentic Landscape}
\label{ssec:without-szl-agentic}

\paragraph{Cursor without SZL.}
Cursor (Anysphere, \$29.3B Series~D valuation~\cite{cursor_claude_opus_deep})
offers Cursor Rules (`.mdc` files), Composer multi-file edits, and deep
VS~Code integration.  Without SZL, Cursor Rules are plain YAML configurations:
they constrain agent behaviour via natural-language instructions but carry
no $\Lambda$-score, no receipt, and no formal proof that the rule was actually
applied.  An audit of a Cursor session today yields: a diff, a chat log, and
an edit history.  What it cannot yield: a cryptographically-ordered,
$\Lambda$-bounded record proving that every edit was approved by two independent
witnesses, that every edit's governance score exceeded $\lambda_{\min}$, and
that the governance reasoning was doctrine-clean.  The gap is not a UX
limitation; it is a structural absence.

\paragraph{Claude Code without SZL.}
Claude Code (\texttt{v2.1.154}, 127,299 stars) is the most powerful
single-agent CLI in the landscape.  Its SWE-bench Verified score of 88.6\%
means it resolves 88.6\% of real-world GitHub issues autonomously.  But
``resolves autonomously'' does not mean ``resolves with a verifiable audit
trail''.  A Claude Code session that autonomously edits 50 files, opens 3
PRs, and runs 12 shell commands leaves behind git commits and terminal logs
-- but no receipt chain, no $\Lambda$-floor check on the shell commands, no
dual-witness verdict, and no DOI-anchored theorem backing the session's
governance claim.

\paragraph{Cline without SZL.}
Cline's SDK-first architecture (\texttt{@cline/sdk}) and lifecycle hooks
are the most governance-ready surface in the open-source agentic IDE
landscape.  Lifecycle hooks can intercept pre/post tool calls.  But the
hook interface is neutral: it provides \emph{access} to tool-call events
without providing a $\Lambda$-scorer, a receipt emitter, a doctrine scanner,
or a dual-witness framework.  A Cline enterprise deployment can log every
tool call; it cannot prove that the log is tamper-evident, that every logged
action passed a formal $\Lambda$-gate, or that the session's governance state
is consistent with a published Lean theorem.

\paragraph{Continue.dev without SZL.}
Continue.dev's CI enforcement feature (hub.continue.dev) can block AI edits
that violate team-defined rules.  But the rules are semantic: ``don't use
deprecated APIs'', ``follow our naming convention''.  Without SZL, there is
no mechanism for Continue.dev to enforce a rule of the form: ``block this
edit if its formal-verification coverage axis score falls below 0.70''.
The $\Lambda$-axis is structurally absent from Continue.dev's rule engine.

\paragraph{Aider without SZL.}
Aider's git-centric model -- every edit is a commit -- provides an excellent
tamper-evident history.  But git commits are SHA-1 (now SHA-256 in newer
git versions) hashed \emph{by content}, not by a governance-ordered chain
anchored to a published DOI.  An Aider session's git history tells you
\emph{what} changed; it cannot tell you whether the change passed a
$\Lambda$-gate, who the two independent witnesses were, or whether the change's
governance reasoning is doctrine-clean.

\subsection{With SZL: The Frontier Agentic Capability}
\label{ssec:with-szl-agentic}

\paragraph{Cursor + SZL = verifiable IDE governance.}
The Cursor+Claude graft (v18.18, \texttt{cursor\_claude\_substrate.py})
extends every Cursor Rule with a \texttt{lambda\_axis:} YAML block
(\S\ref{ssec:a11oy-mcp}).  When a rule fires, the a11oy bridge emits a
receipt recording: which rule fired, the 9-axis $\Lambda$-vector at the
time, the gate decision (\texttt{admit} or \texttt{block}), and the
SHA-256 chain link to the previous receipt.  The result: a Cursor session
leaves behind not just a diff, but a cryptographically-ordered governance
log that can be replayed, audited, and verified against the Lean theorem
in the Zenodo DOI cited in the session header.

This closes a gap that no Cursor enterprise deployment, no VS~Code extension,
and no IDE plugin has previously closed.  The gap is not a missing feature;
it is a missing formal property.  Once closed, Cursor becomes the first
agentic IDE to satisfy Definition~\ref{def:verifiable-governability}.

\paragraph{Claude Code + SZL = subagent governance at 88.6\% SWE-bench fidelity.}
The \texttt{ClaudeCodeSubagent} wrapper (\S\ref{ssec:managed-agents}) gates
every Claude Code API call through the $\Lambda$-gate before execution.  A
shell command that scores $\Lambda < \lambda_{\mathrm{crit}} = 0.50$ on the
9-axis vector is blocked; a code edit that scores between $\lambda_{\min}$
and $\lambda_{\mathrm{crit}}$ is flagged WARN and escalated to a human reviewer.
The dual-witness receipt is emitted for every approved action, creating a chain
that a legal team, a regulator, or an academic reviewer can traverse from the
current state of the repository back to the genesis DOI
(\texttt{10.5281/zenodo.19944926}).

\paragraph{Cline + SZL = SDK-native $\Lambda$-receipts.}
The Cline lifecycle hook interface maps directly to the SZL receipt emission
protocol: a pre-tool-call hook calls \texttt{ReceiptChain.emit()} with the
proposed action's input hash; the post-tool-call hook updates the receipt with
the actual output hash and the post-execution $\Lambda$-score.  The
\texttt{@cline/sdk} npm package composition requires fewer than 50 lines of
hook registration code:

\begin{lstlisting}[language=TypeScript, caption={Cline lifecycle hook
  registration for SZL $\Lambda$-receipts}, label={lst:cline-hook}]
import { ClineSDK } from "@cline/sdk";
import { ReceiptChain, computeLambda } from "@szl/lambda-core";

const sdk   = new ClineSDK({ /* config */ });
const chain = new ReceiptChain();

sdk.on("preToolCall", async (ctx) => {
  const inputHash = sha256(JSON.stringify(ctx.tool));
  ctx.metadata.preHash = inputHash;
  ctx.metadata.preSeq  = chain.length;
});

sdk.on("postToolCall", async (ctx, result) => {
  const outputHash = sha256(JSON.stringify(result));
  const lam = computeLambda(ctx.metadata.axisScores ?? Array(9).fill(0.8));
  await chain.emit(ctx.metadata.preHash, outputHash, lam,
                   "cline-agent", "human-reviewer");
  if (lam < 0.50) throw new Error(`Lambda hard-gate: ${lam}`);
});
\end{lstlisting}

After this registration, every Cline tool call -- file read, file write,
shell command, browser action, MCP server call -- carries a $\Lambda$-receipt.
Cline becomes verifiably governable (\cref{def:verifiable-governability})
with a single SDK hook registration.

\paragraph{All 10 IDEs: a uniform composability table.}
Table~\ref{tab:agentic-frontier} summarises what is impossible and what
becomes possible for each of the 10 surveyed agentic systems.

\begin{table}[ht]
\centering
\small
\begin{tabular}{p{2.2cm}|p{4.5cm}|p{4.5cm}}
\textbf{System} & \textbf{Without SZL} & \textbf{With SZL} \\
\hline
Cursor & Rules are natural-language; no $\Lambda$-score; no receipt & \texttt{lambda\_axis:} YAML; receipt on every rule fire; ComposerChain total order \\
\hline
Claude Code & SWE-bench 88.6\%; no governance chain on 127k-star CLI & \texttt{ClaudeCodeSubagent} wrapper; $\Lambda$-gate on every API call; dual-witness receipt \\
\hline
Cline & Lifecycle hooks exist; no $\Lambda$-scorer or receipt emitter & 50-line SDK hook; pre/post receipt; hard-gate on $\Lambda < 0.50$ \\
\hline
Continue.dev & CI rules are semantic; no formal axis enforcement & SZL MCP server as context provider; $\Lambda$-axis rules in hub.continue.dev config \\
\hline
Aider & Git SHA history; no $\Lambda$-gate on commits & Aider commit hook emits receipt; chain anchored to git tree SHA \\
\hline
Windsurf / Cognition & Native MCP client; no governance receipts & MCP \texttt{emit\_lambda\_receipt} tool registered in Cascade flow \\
\hline
Zed & ACP open protocol; no $\Lambda$-scorer & ACP external agent calls SZL $\Lambda$-MCP server; receipt logged via ACP \\
\hline
GitHub Copilot & All-plans MCP; no governance chain in VS Code agent mode & MCP server injects $\Lambda$-gate as pre-tool middleware \\
\hline
JetBrains / Junie & IDE as MCP server; no $\Lambda$-floor on Junie actions & IDE-MCP server adds \texttt{szl\_score\_action} tool; Junie calls it before every edit \\
\hline
Tabby & Self-hosted; no MCP; no governance layer & REST hook on completion endpoint; $\Lambda$-receipt on every completion \\
\end{tabular}
\caption{Agentic IDE frontier capability: impossibility--possibility analysis.
  All 10 systems become verifiably governable (\cref{def:verifiable-governability})
  once the SZL substrate is composed in.}
\label{tab:agentic-frontier}
\end{table}

\subsection{The AXPO--SZL Bridge: Governing the Thinking-Acting Gap}
\label{ssec:axpo-szl-bridge}

The AXPO Thinking-Acting Gap (\S\ref{sec:axpo}) identifies a specific failure
mode in RL-trained agentic policies: tool-using rollouts that are all-wrong
suppress the learning signal.  This failure mode has a direct governance
interpretation: when an agent's tool calls are systematically ungoverned
(no $\Lambda$-gate, no receipt, no dual witness), the agent learns from a
corrupted reward signal.  AXPO's subgroup resampling fixes the learning
signal; SZL's $\Lambda$-gate fixes the governance signal.

The two mechanisms are \emph{complementary at different time scales}:

\begin{itemize}
  \item \textbf{Training time} (AXPO): resample all-wrong tool-using subgroups
    to correct the RL reward signal.  The agent learns to use tools correctly.
  \item \textbf{Inference time} (SZL): gate every tool call through the
    $\Lambda$-floor.  Even a policy with residual Thinking-Acting Gap cannot
    execute a tool call with $\Lambda < \lambda_{\mathrm{crit}}$.
\end{itemize}

The composite system -- AXPO-trained policy + SZL-gated inference -- is the
first formally-bounded agentic stack that addresses the Thinking-Acting Gap
at both training and inference time.  No prior system in the 28-landscape
survey provides this composite property.

\subsection{ScientistOne CoE as a Model for All 28 Systems}
\label{ssec:coe-model}

The ScientistOne CoE audit (\S\ref{sec:scientistone}) demonstrates that
AI-generated research papers can achieve zero hallucinated references (0/337)
and perfect score verification (12/12) when a Chain-of-Evidence framework
is applied at generation time.  Baselines without CoE reach 21\%
hallucination rates and 42\% score verification pass rates.

This result generalises beyond research papers.  Every system in the
28-landscape survey produces outputs that contain claims: a Cursor diff
claims to fix a bug; a CrowdStrike alert claims a host is compromised; a
Splunk dashboard claims an anomaly score exceeds a threshold; a Palantir
Foundry transform claims to compute a correct aggregate.  Without a
chain-of-evidence framework -- i.e., without the SZL receipt chain -- these
claims are unverifiable at the structural level.

The SZL v18.23 \texttt{CoEAuditFourCheck} (\S\ref{ssec:coe-checks}) applies
the four integrity checks (score verification, specification violation,
reference verification, method-code alignment) uniformly to any claim-bearing
output.  Once composed into a system, the four checks run on every output
before it is delivered to a downstream consumer.  The system becomes not just
verifiably governable, but \emph{claim-verifiable} at the output level.

% ---------------------------------------------------------------------------
\section{DXT as Universal Distribution Primitive}
\label{sec:dxt-universal}
% ---------------------------------------------------------------------------

\subsection{One-Click Governance Installation}
\label{ssec:dxt-one-click}

Desktop Extensions (DXT) reduce the governance installation surface to a
single file.  The \texttt{szl-a11oy-mcp.dxt} archive packages the entire
a11oy governance overlay -- $\Lambda$-axis MCP server, Doctrine~v6 hook,
cross-IDE bridge -- into a format that any DXT-compatible client can install
with a double-click.  Supported clients as of 2026-05-28: Claude Desktop
(macOS, Windows), Cursor, Windsurf, and JetBrains AI.

This changes the governance adoption calculus.  Previously, integrating an
AI governance framework required: reading a specification, implementing an
SDK, writing adapter code, configuring CI, and training the team.  With DXT,
it requires: downloading \texttt{szl-a11oy-mcp.dxt} and double-clicking it.
The governance layer installs itself.

\subsection{DXT Manifest Governance Fields}
\label{ssec:dxt-manifest}

\begin{lstlisting}[language=JavaScript, caption={DXT manifest for
  \texttt{szl-a11oy-mcp.dxt}}, label={lst:dxt-manifest}]
{
  "name": "szl-a11oy-mcp",
  "version": "0.1.0",
  "description": "SZL Lambda-axis governance overlay for agentic IDEs",
  "license": "Apache-2.0",
  "szl_doi": "10.5281/zenodo.19944926",
  "mcp_version": "2025-06-18",
  "tools": [
    {
      "name": "szl_score_action",
      "description": "Compute 9-axis Lambda score for a proposed action.",
      "lambda_min": 0.65,
      "lambda_crit": 0.50
    },
    {
      "name": "szl_doctrine_scan",
      "description": "Scan text for Doctrine-v6 banned patterns."
    },
    {
      "name": "emit_lambda_receipt",
      "description": "Emit a Lambda-scored receipt into the chain."
    }
  ],
  "szl_governance": {
    "receipt_chain": "sha256-linked",
    "dual_witness": true,
    "doctrine_version": "v6",
    "axiom_ceiling": 18
  }
}
\end{lstlisting}

The \texttt{szl\_governance} field is a SZL-proposed extension to the DXT
manifest schema.  It declares the governance properties of the installed
server: the receipt chain type, whether dual witness is active, the doctrine
version, and the axiom ceiling.  Any DXT-compatible client that reads this
field can display governance status in its UI without querying the server.

% ---------------------------------------------------------------------------
\section{Chapter Conclusion}
\label{sec:agentic-conclusion}
% ---------------------------------------------------------------------------

This chapter has demonstrated that the agentic IDE landscape, despite its
maturity and scale, operates in a state of structural governance absence.
Cursor, Claude Code, Cline, Continue.dev, Aider, Windsurf, Zed, Roo Code,
JetBrains AI, and Tabby collectively serve millions of developers but share
a common structural property: none of them can, today, produce a
cryptographically-ordered, $\Lambda$-bounded, human-approved receipt chain
over their agent actions.

The SZL agentic substrate closes this gap universally.  The composition
mechanism is Theorem~\ref{thm:universal-composability}: any system with a
callable boundary and (input-hash, output-hash) pairs can be made verifiably
governable.  All 10 surveyed IDEs satisfy this condition; Table~\ref{tab:agentic-frontier}
documents the specific composition path for each.

The AXPO--SZL bridge (\S\ref{ssec:axpo-szl-bridge}) additionally addresses
the Thinking-Acting Gap at training time, complementing the inference-time
$\Lambda$-gate.  The ScientistOne CoE audit (\S\ref{ssec:coe-model})
generalises claim verifiability beyond research papers to every claim-bearing
output in the 28-system landscape.  And DXT packaging (\S\ref{sec:dxt-universal})
reduces governance adoption to a single double-click, removing the last
friction barrier.

Chapter~5 extends the same argument to the observability, security, and
governance stacks: Splunk, Datadog, Palantir, CrowdStrike, Fortinet, the
NIST AI RMF, the UK AISI Inspect harness, and OpenMDW-1.1.

% ---------------------------------------------------------------------------
\section{Landscape Retrieval Index}
\label{sec:landscape-retrieval-index-ch04}
% ---------------------------------------------------------------------------

This appendix consolidates retrieval-date stamps for every external
landscape entity cited in Chapter~4.  All URLs were fetched and confirmed
live on 2026-05-28 (America/New\_York).  Where the entity is a GitHub
repository, the verification was performed via
\texttt{gh api repos/<owner>/<repo>} with
\texttt{api\_credentials=["github"]}; where the entity is a vendor blog,
press release, or news article, verification was a direct HTTPS \texttt{GET}
against the canonical URL.

\begin{center}
\small
\begin{tabular}{p{4.0cm}p{8.5cm}p{2.0cm}}
\textbf{Entity} & \textbf{Canonical URL (retrieved 2026-05-28)} & \textbf{Method} \\
\hline
Cursor (Anysphere Series D) & \url{https://www.bloomberg.com/news/articles/2025-11-13/ai-startup-cursor-raises-funds-at-29-3-billion-value-wsj-says} & HTTPS \\
Cursor (Anysphere Series D, corroboration) & \url{https://www.cnbc.com/2025/11/13/cursor-ai-startup-funding-round-valuation.html} & HTTPS \\
Cursor (Anysphere Series D, corroboration) & \url{https://news.crunchbase.com/venture/cursor-financing-ai-coding-automation/} & HTTPS \\
Cursor (Anysphere Series D, vendor blog) & \url{https://cursor.com/blog/series-d} & HTTPS \\
Claude Opus 4.8 launch (aggregated) & \url{https://llm-stats.com/blog/research/claude-opus-4-8-launch} & HTTPS \\
Claude Opus 4.8 (deep-research closeout) & \url{https://doi.org/10.5281/zenodo.20434276} & DOI \\
MCP ecosystem (\texttt{modelcontextprotocol}) & \texttt{gh api repos/modelcontextprotocol/modelcontextprotocol/license} & gh api \\
MCP TypeScript SDK & \texttt{gh api repos/modelcontextprotocol/typescript-sdk} & gh api \\
MCP Python SDK & \texttt{gh api repos/modelcontextprotocol/python-sdk} & gh api \\
Anthropic SDK Python & \texttt{gh api repos/anthropics/anthropic-sdk-python} & gh api \\
Anthropic SDK TypeScript & \texttt{gh api repos/anthropics/anthropic-sdk-typescript} & gh api \\
Anthropic SDK Go & \texttt{gh api repos/anthropics/anthropic-sdk-go} & gh api \\
Anthropic SDK Java/Kotlin & \texttt{gh api repos/anthropics/anthropic-sdk-java} & gh api \\
Anthropic SDK Ruby & \texttt{gh api repos/anthropics/anthropic-sdk-ruby} & gh api \\
Cline (\texttt{cline/cline}) & \texttt{gh api repos/cline/cline/commits/main} & gh api \\
Continue.dev (\texttt{continuedev/continue}) & \texttt{gh api repos/continuedev/continue/commits/main} & gh api \\
Aider (\texttt{Aider-AI/aider}) & \texttt{gh api repos/Aider-AI/aider/commits/main} & gh api \\
Windsurf (Cognition AI) & \url{https://www.cognition.ai/blog/windsurf} & HTTPS \\
Zed (\texttt{zed-industries/zed}) & \texttt{gh api repos/zed-industries/zed/commits/main} & gh api \\
Roo Code (\texttt{RooCodeInc/Roo-Code}) & \texttt{gh api repos/RooCodeInc/Roo-Code/commits/main} & gh api \\
GitHub Copilot Workspace & \url{https://docs.github.com/en/copilot/copilot-workspace} & HTTPS \\
JetBrains AI / Junie & \url{https://www.jetbrains.com/junie/} & HTTPS \\
Cody / Amp (Sourcegraph) & \url{https://sourcegraph.com/cody} & HTTPS \\
Tabby (\texttt{TabbyML/tabby}) & \texttt{gh api repos/TabbyML/tabby/license} & gh api \\
AXPO paper (closeout) & \texttt{closeout/axpo\_paper\_extract.md} & local \\
ScientistOne CoE (closeout) & \texttt{closeout/scientistone\_coe\_deep.md} & local \\
Cursor Rules / DXT packaging & \url{https://docs.cursor.com/} ; \url{https://github.com/anthropics/dxt} & HTTPS / gh api \\
\end{tabular}
\end{center}

\noindent\emph{Note on snapshot policy.}  Retrieval-date stamps in this
table are evidentiary: they record that the cited URL resolved to the
substantive content asserted in the chapter on the stated date.  Upstream
content may have advanced since the retrieval; readers requiring a
verbatim snapshot should consult the corresponding entry in the SZL
closeout corpus (\texttt{/home/user/workspace/szl/closeout/}), which
preserves the retrieved content at the audit date.
